%% kapa: exam=1 round=37 year=2026 period=1 form=A subject=economics subject_ko=경제학원론 seq=11 number=51
%% source: Q-Net 감정평가사 기출문제 (cst00309) artlSeq=5258924 file=raw/1008/5258924/2026년 제37회 감정평가사 1차 1교시.pdf
%% conversion: kordoc 4.24.51 -> tools/gen1.py (auto); review=proofread-sonnet; flags=none
\begin{question}{51}
A, B 두 나라에 각각 다섯 사람씩 살고 있고 A국과 B국에 사는 사람들의 소득은 각각 (2, 2, 4, 4, 8), (2, 4, 4, 4, 6)이라고 하자. 다음 설명 중 옳지 않은 것은?
\begin{choices}
  \item 지니계수상으로나 오분위분배율로나 B국이 더 균등하다.
  \item A국의 십분위분배율은 0.5이고 B국의 오분위분배율은 3이다.
  \item B국의 십분위분배율은 1이고 A국의 오분위분배율은 4이다.
  \item 오분위분배율로 보면 B국이 더 균등하며, 십분위분배율은 A국이 더 불균등하다.
  \item 지니계수상으로 A국이 더 균등하며, 십분위분배율로 보면 B국이 더 균등하다.
\end{choices}
\end{question}
\answer{5}
